\EMhold

\EMsection{The Set of Complex Numbers}{The Set of Complex Numbers}

\begin{EMdefinition}{Complex numbers}

The set of complex numbers generalizes the real numbers on the real line to a plane of numbers having two components, a real one and an imaginary one.
\EMdisp{\mathbb{R}\to\mathbb{C}}

\end{EMdefinition}

\begin{EMunit}

\EMdisp{\mathbb{C}\EMbr =\{(x,y)\mid x,y\in\mathbb{R}\},\qquad\EMbr  z\EMbr =(x,y),\qquad\EMbr  x\EMbr =\operatorname{Re}\{z\},\qquad\EMbr  y\EMbr =\operatorname{Im}\{z\}}

\EMfigure{m08-complex-plane}{The complex number \EMim{z=(x,y)} in the complex plane; its modulus
\EMim{|z|} and argument \EMim{\theta}}

\end{EMunit}

\begin{EMunit}

Magnitude or norm of a complex number:
\EMdisp{|z|\EMbr =\sqrt{x^2+y^2}}

Phase or argument:
\EMdisp{\theta\EMbr =\measuredangle z\EMbr =\tan^{-1}\!\left(\frac{y}{x}\right),\qquad\EMbr  z\EMbr =|z|\measuredangle\theta}
\EMdisp{x\EMbr =|z|\cos\theta,\qquad\EMbr  y\EMbr =|z|\sin\theta}

\EMdisp{z\EMbr =x\times(1,0)+y\times(0,1)\EMbr =x+iy}

\end{EMunit}

\begin{EMimportant}{}

\EMdisp{i^2\EMbr =-1}

\end{EMimportant}

\EMhold

\EMsection{The Four Basic Operations}{The Four Basic Operations}

\begin{EMunit}

Let \EMim{z_1=(a_1,b_1)=a_1+ib_1} and \EMim{z_2=(a_2,b_2)=a_2+ib_2}.

\textbf{(I) Addition:}
\EMdisp{z_1+z_2\EMbr =(a_1+a_2,\,b_1+b_2)\EMbr =a_1+a_2+i(b_1+b_2)}

\end{EMunit}

\begin{EMunit}

\textbf{(II) Subtraction:}
\EMdisp{z_1-z_2\EMbr =(a_1-a_2,\,b_1-b_2)\EMbr =a_1-a_2+i(b_1-b_2)}

\textbf{(III) Multiplication:}
\EMdisp{\begin{aligned}
z_1\times z_2&=(a_1+ib_1)\times(a_2+ib_2)\\
&=a_1\times a_2+a_1\times ib_2+ib_1\times a_2+ib_1\times ib_2\\
&=(a_1a_2-b_1b_2)+i(a_1b_2+b_1a_2)
\end{aligned}}

\end{EMunit}

\begin{EMunit}

Multiplication by a real number:
\EMdisp{\forall\alpha\in\mathbb{R}:\quad\EMbr  \alpha z_1\EMbr =\alpha a_1+i\alpha b_1}

\textbf{(IV) Division:}
\EMdisp{\begin{aligned}
z_1\div z_2&=(a_1+ib_1)\div(a_2+ib_2)=\frac{a_1+ib_1}{a_2+ib_2}=\frac{(a_1+ib_1)(a_2-ib_2)}{(a_2+ib_2)(a_2-ib_2)}\\
&=\frac{(a_1a_2+b_1b_2)+i(b_1a_2-a_1b_2)}{a_2^2+b_2^2}
=\frac{a_1a_2+b_1b_2}{a_2^2+b_2^2}+i\,\frac{b_1a_2-a_1b_2}{a_2^2+b_2^2}
\end{aligned}}

\end{EMunit}

\EMhold

\EMsection{Properties of Addition and Multiplication}{Properties of Addition and Multiplication}

\begin{EMtheorem}{The complex numbers form a field}

The set of complex numbers, together with addition and multiplication, forms a field that cannot be ordered.

\end{EMtheorem}

\begin{EMunit}

For \EMim{z_1,z_2,z_3\in\mathbb{C}}:

\EMdisp{z_1+z_2\in\mathbb{C},\qquad\EMbr  z_1\times z_2\in\mathbb{C}}
\EMdisp{z_1+z_2\EMbr =z_2+z_1,\qquad\EMbr  z_1\times z_2\EMbr =z_2\times z_1}
\EMdisp{(z_1+z_2)+z_3\EMbr =z_1+(z_2+z_3),\qquad\EMbr  (z_1z_2)z_3\EMbr =z_1(z_2z_3)}
\EMdisp{z_1(z_2+z_3)\EMbr =z_1z_2+z_1z_3}
\EMdisp{z_1+(0,0)\EMbr =z_1+0\EMbr =z_1,\qquad\EMbr  z_1+(-z_1)\EMbr =(0,0)\EMbr =0}
\EMdisp{-z_1\EMbr =-(a_1,b_1)\EMbr =(-a_1,-b_1)\EMbr =-a_1-ib_1}
\EMdisp{z_1\times(1,0)\EMbr =z_1\times1\EMbr =z_1,\qquad\EMbr  z_1\times z_1^{-1}\EMbr =(1,0)\EMbr =1}
\EMdisp{z_1^{-1}\EMbr =\left(\frac{a_1}{a_1^2+b_1^2},\frac{-b_1}{a_1^2+b_1^2}\right)\EMbr =\frac{a_1}{a_1^2+b_1^2}-i\,\frac{b_1}{a_1^2+b_1^2}}

\end{EMunit}

\begin{EMremark}{No ordering}

An order relation \EMim{z_1<z_2} cannot be defined on the set of complex numbers.

\end{EMremark}

\EMhold

\EMsection{The Complex Conjugate}{The Complex Conjugate}

\begin{EMdefinition}{Complex conjugate of a number}

\EMdisp{z\EMbr =x+iy\quad\EMbr \EMbr \Longrightarrow\quad\EMbr  \bar z\triangleq x-iy}

\end{EMdefinition}

\begin{EMexample}{}

Plot the conjugate of the complex number \EMim{z=5+2i} in the complex plane.

\end{EMexample}

\begin{EMunit}

\EMfigure{m08-conjugate}{The complex number \EMim{z=5+2i} and its conjugate \EMim{\bar z=5-2i}}

\end{EMunit}

\begin{EMexercise}{}

Prove the following statements.
\EMdisp{\operatorname{Re}\{z\}\EMbr =\frac12\,(z+\bar z),\qquad\EMbr  \operatorname{Im}\{z\}\EMbr =\frac{1}{2i}\,(z-\bar z),\qquad\EMbr  |z|^2\EMbr =z\cdot\bar z}
\EMdisp{\overline{z_1+z_2}\EMbr =\bar z_1+\bar z_2,\qquad\EMbr  \overline{z_1\cdot z_2}\EMbr =\bar z_1\cdot\bar z_2,\qquad\EMbr  \overline{\left(\frac{z_1}{z_2}\right)}\EMbr =\frac{\bar z_1}{\bar z_2}}
\EMdisp{|z|\EMbr =|\bar z|,\qquad\EMbr  |\operatorname{Re}\{z\}|\le|z|,\qquad\EMbr  |\operatorname{Im}\{z\}|\le|z|}

\end{EMexercise}

\EMhold

\EMsection{The Triangle Inequality}{The Triangle Inequality}

\begin{EMunit}

\EMfigure{m08-triangle}{The triangle inequality: a triangle with vertices \EMim{0}, \EMim{z_1}
and \EMim{z_1+z_2}}

\end{EMunit}

\begin{EMexercise}{}

\begin{itemize}
\tightlist
\item
  Prove the triangle inequality for any two complex numbers: \EMim{|z_1+z_2|\le|z_1|+|z_2|}.
\item
  Prove the general form of the triangle inequality for \EMim{n} arbitrary complex numbers:
  \EMdisp{|z_1+z_2+\cdots+z_n|\le|z_1|+|z_2|+\cdots+|z_n|}
\item
  Prove that \EMim{|z_1+z_2|\ge|z_1|-|z_2|}.
\end{itemize}

\end{EMexercise}

\EMhold

\EMsection{The Polar Form of Complex Numbers}{The Polar Form of Complex Numbers}

\begin{EMunit}

\EMdisp{z\EMbr =x+iy,\qquad\EMbr  \left.\begin{aligned}x&=|z|\cos\theta\\ y&=|z|\sin\theta\end{aligned}\right\}\;\EMbr \Longrightarrow\; z\EMbr =|z|\,(\cos\theta+i\sin\theta)}

\EMdisp{\theta\EMbr =\operatorname{Arg}(z)\EMbr =\tan^{-1}\frac yx,\qquad\EMbr  -\pi<\theta\le\pi}
\EMdisp{\arg(z)\EMbr =\operatorname{Arg}(z)+2k\pi,\qquad\EMbr  k\in\mathbb{Z}}

\end{EMunit}

\EMhold

\EMsection{Multiplying Two Complex Numbers in Polar Form}{Multiplying Two Complex Numbers in Polar Form}

\begin{EMunit}

\EMdisp{z_1\EMbr =|z_1|(\cos\theta_1+i\sin\theta_1),\qquad\EMbr  z_2\EMbr =|z_2|(\cos\theta_2+i\sin\theta_2)}

\EMdisp{\begin{aligned}
z_1\cdot z_2&=|z_1||z_2|(\cos\theta_1+i\sin\theta_1)(\cos\theta_2+i\sin\theta_2)\\
&=|z_1||z_2|\big[(\cos\theta_1\cos\theta_2-\sin\theta_1\sin\theta_2)\\
&\qquad\qquad+\,i(\sin\theta_1\cos\theta_2+\cos\theta_1\sin\theta_2)\big]
\end{aligned}}

\end{EMunit}

\begin{EMimportant}{}

\EMdisp{z_1\cdot z_2\EMbr =|z_1||z_2|\big[\cos(\theta_1+\theta_2)+i\sin(\theta_1+\theta_2)\big]}

\end{EMimportant}

\begin{EMremark}{Consequence 1}

When two complex numbers are multiplied, their magnitudes are multiplied and their phases are added.

\end{EMremark}

\begin{EMremark}{Consequence 2}

If a complex number is raised to a natural power, its magnitude is raised to that power and its argument (phase) is multiplied by it (prove by induction). In other words:
\EMdisp{z^k\EMbr =|z|^k(\cos k\theta+i\sin k\theta)}

\end{EMremark}

\begin{EMunit}

\EMdisp{z^k\EMbr =\big[|z|(\cos\theta+i\sin\theta)\big]^k\EMbr =|z|^k(\cos\theta+i\sin\theta)^k\EMbr =|z|^k(\cos k\theta+i\sin k\theta)}

\end{EMunit}

\begin{EMimportant}{De Moivre's formula}

\EMdisp{(\cos\theta+i\sin\theta)^k\EMbr =\cos k\theta+i\sin k\theta}

\end{EMimportant}

\begin{EMexercise}{}

Using De Moivre's formula, express \EMim{\cos3\theta} and \EMim{\cos4\theta} in terms of \EMim{\cos\theta} and \EMim{\sin\theta}.

\end{EMexercise}

\EMhold

\EMsection{The \EMim{m}-th Root of a Complex Number}{The m-th Root of a Complex Number}

\begin{EMunit}

\EMdisp{w\EMbr =\sqrt[m]{z}\EMbr =z^{1/m},\qquad\EMbr  z\EMbr =|z|(\cos\theta+i\sin\theta)}

\EMdisp{z\EMbr =w^m,\qquad\EMbr  w\EMbr =|w|(\cos\varphi+i\sin\varphi)}

\EMdisp{z\EMbr =|z|(\cos\theta+i\sin\theta)\EMbr =|w|^m(\cos m\varphi+i\sin m\varphi)}

\EMdisp{|w|^m\EMbr =|z|\;\EMbr \Longrightarrow\;|w|\EMbr =\sqrt[m]{|z|}}

\EMdisp{m\varphi\EMbr =\theta+2k\pi\;\EMbr \Longrightarrow\;\varphi\EMbr =\frac{\theta+2k\pi}{m},\qquad\EMbr  k\EMbr =0,1,2,\dots,m-1}

\end{EMunit}

\begin{EMimportant}{}

\EMdisp{w\EMbr =\sqrt[m]{z}\EMbr =\sqrt[m]{|z|}\left[\cos\!\left(\frac{\theta+2k\pi}{m}\right)+i\sin\!\left(\frac{\theta+2k\pi}{m}\right)\right],\qquad\EMbr  k\EMbr =0,1,2,\dots,m-1}

\end{EMimportant}

\begin{EMexample}{}

What are the roots of the equation \EMim{z^n=1}?

\EMdisp{z\EMbr =\sqrt[n]{1},\qquad\EMbr  1\EMbr =1\measuredangle0\;\EMbr \Longrightarrow\; z\EMbr =\cos\!\left(\frac{2k\pi}{n}\right)+i\sin\!\left(\frac{2k\pi}{n}\right),\qquad\EMbr  k\EMbr =0,1,\dots,n-1}
Defining \EMim{\omega=1\measuredangle\dfrac{2\pi}{n}}, the roots are
\EMdisp{z\EMbr =1,\ \omega,\ \omega^2,\ \dots,\ \omega^{n-1}}

\end{EMexample}

\begin{EMunit}

\EMfigure{m08-roots-of-unity}{The roots \EMim{z=\sqrt[3]{1}}, \EMim{z=\sqrt[4]{1}} and
\EMim{z=\sqrt[5]{1}} on the unit circle}

\end{EMunit}

\begin{EMexercise}{}

\begin{itemize}
\tightlist
\item
  Find the roots of the equation \EMim{z^6+6+8i=0} and plot them in the complex plane.
\item
  Verify the identity \EMim{1+z+z^2+\cdots+z^n=\dfrac{1-z^{n+1}}{1-z}} and then derive the following identity (Lagrange's identity).
  \EMdisp{1+\cos\theta+\cos2\theta+\cdots+\cos n\theta\EMbr =\frac12+\frac{\sin\!\left[\left(n+\frac12\right)\theta\right]}{2\sin\frac\theta2}}
\end{itemize}

\end{EMexercise}
